\documentclass[11pt,a4paper]{article}
\usepackage[T1]{fontenc}
\usepackage{lmodern}
\usepackage[margin=25mm]{geometry}
\usepackage{amsmath,amssymb,amsthm,booktabs,array}
\usepackage{microtype}
\usepackage[hidelinks]{hyperref}
\hypersetup{
  pdftitle={Exact bounded fibres and the cost of primitive Wronskians},
  pdfauthor={Oleksiy Babanskyy},
  pdfsubject={Arithmetic derivatives, bounded integer fibres and primitive Wronskian costs},
  pdfkeywords={arithmetic derivatives, integer lattices, bounded fibres, Wronskians, exact certificates}
}
\pagestyle{plain}
\newtheorem{theorem}{Theorem}[section]
\newtheorem{proposition}[theorem]{Proposition}
\newtheorem{lemma}[theorem]{Lemma}
\newtheorem{corollary}[theorem]{Corollary}
\theoremstyle{definition}\newtheorem{question}[theorem]{Question}
\theoremstyle{remark}\newtheorem{remark}[theorem]{Remark}
\newcommand{\Z}{\mathbb Z}\newcommand{\R}{\mathbb R}
\newcommand{\Dgen}{\mathcal D}\newcommand{\mpr}{m_{\mathrm{prim}}}
\newcommand{\norm}[1]{\left\lVert#1\right\rVert_\infty}
\newcommand{\hh}{\mathfrak h}\newcommand{\ee}{\mathrm e}
\DeclareMathOperator{\rad}{rad}\DeclareMathOperator{\supp}{supp}
\DeclareMathOperator{\lcm}{lcm}\DeclareMathOperator{\ord}{ord}
\title{Exact bounded fibres and the cost of primitive Wronskians}
\author{Oleksiy Babanskyy}
\date{}
\begin{document}\maketitle
\begin{abstract}
An arithmetic derivative additive on a fixed coprime equation $a+b=c$
may have a nonzero Wronskian without representing a generator of its
integer image. We compare the least maximum norms for these two tasks
in the original prime coordinates. For an odd prime $q$ such that
$q+2$ has a prime divisor of exponent one, the nonzero minimum for
$2+q=q+2$ is one, while the primitive minimum is exactly
$(q-\Dgen)/2$ when $q+2$ is composite and $(q+1)/2$ when it is prime.
The image generator is explicit, and the excluded prime indices have
relative density zero. We also give exact bounded-fibre formulas and
extend the analysis to the radical-small triples $1+(x^N-1)=x^N$,
with $N=4\rad_{\rm odd}(x+1)$. A uniform comparison identifies their
minimum ratio with a simultaneous modular-return problem on a
positive-density set of bases. In this family the number of large
active coordinates grows in density; at feasible subquadratic budgets,
fixed-order frequency interactions and polynomially many explicit
Fourier terms cannot certify a return outside a set of density zero. Exact certificates give
the two minima for $2+3^{47}$ and for the bases $x=14,44$.
Unbounded separation in the fixed-power family and in the radical-small
regime remains open. No consequence for the $abc$ conjecture is claimed.
\end{abstract}

\section{The two minimum problems}\label{sec:intro}
Let $1\le a\le b$, $a+b=c>2$, and $\gcd(a,b)=1$. The prime supports
of $a,b,c$ are pairwise disjoint. In the model of Pasten
\cite[Sections 2.2 and 2.4]{Pasten}, a derivative supported on
$S=\supp(abc)$ is specified by integers $z=(z_p)_{p\in S}$. Set
\begin{align}
 D_n(z)&=\sum_{p\mid n}\frac{n v_p(n)}p z_p,&
 F(z)&=D_a(z)+D_b(z)-D_c(z),\label{eq:model}\\
 W(z)&=aD_b(z)-bD_a(z),&
 \Lambda&=\ker_{\Z}F,\quad \Lambda_0=\Lambda\cap\ker_{\Z}W.\nonumber
\end{align}
Empty sums are zero. All norms are $\norm{z}=\max_{p\in S}|z_p|$
in these original prime coordinates. The independent rows $F,W$ have
coefficient ratios $-b,a,0$ on their nonempty blocks, of which there
are at least two. Consequently $W(\Lambda)=\Dgen\Z$ for some
$\Dgen>0$. Define
\begin{equation}\label{eq:minima}
 m(a,b)=\min_{\substack{z\in\Lambda\\W(z)\ne0}}\norm{z},\qquad
 \mpr(a,b)=\min_{\substack{z\in\Lambda\\W(z)=\Dgen}}\norm{z}.
\end{equation}
Both minima exist. In this paper \emph{primitive} means a generator
of the image, equivalently of $\Lambda/\Lambda_0$. It does not mean
merely that the coordinates of $z$ have gcd one. Changing the sign of
the specified generator does not change its minimum.

\begin{theorem}[Exact primitive cost for almost all prime indices]\label{thm:main}
Let $q\ge3$ be prime and suppose $c=q+2$ is not powerful, that is,
some prime divisor of $c$ occurs with exponent one. Put
$R_c=\prod_{p\mid c,\ p\nmid v_p(c)}p$ and $D=c/R_c$. Then
\begin{equation}\label{eq:main}
 \Dgen=D,\qquad m(2,q)=1,\qquad
 \mpr(2,q)=
 \begin{cases}(q-D)/2,&c\text{ composite},\\(q+1)/2,&c\text{ prime}.
 \end{cases}
\end{equation}
In the composite case every minimizer for $W=D$ has
\begin{equation}\label{eq:main-opts}
 z_2=-1,\quad z_q=-(q-D)/2,\quad D_c(z)=(D-c)/2,
 \quad |z_p|\le(q-D)/2\quad(p\mid c),
\end{equation}
and every vector satisfying these conditions is a minimizer. In the
prime case the unique positive-generator minimizer is
$(-1,-(q-1)/2,-(q+1)/2)$ in coordinates $(2,q,c)$.
Among prime indices $q\le Q$ fewer than $3\sqrt{Q+2}$ are excluded.
Thus, on a set of relative density one among primes,
\begin{equation}\label{eq:prime-gap}
 m=1,\qquad (q-1)/3\le\mpr\le(q+1)/2.
\end{equation}
\end{theorem}

The norm-one pattern $z_2=-1,z_q=1$, with the $c$-block zero,
has $F=0,W=c$. The substantive part of the theorem is the bounded
lift of the primitive value, not the mere computation of its image.
The elementary case $v_3(q+2)=1$ is contained in its composite branch.
All these triples satisfy $\rad(2q(q+2))\ge6q>q+2$.

The extension to radical-small triples has a different conclusion.
For $N=4\rad_{\rm odd}(x+1)$ and $b=x^N-1$, we construct an exact
simultaneous-return function $h_x(u)$, with $G_x=h_x(1)$ and
$\tau_x=\min_{u\ge1}h_x(u)$. On a positive-density set of bases
there is a constant $C_0$, independent of $x$, such that
\begin{equation}\label{eq:intro-compare}
 \frac{G_x}{C_0\tau_x}\le\frac{\mpr(1,b)}{m(1,b)}
 \le\frac{C_0G_x}{\tau_x}.
\end{equation}
Theorem~\ref{thm:comparison} gives the effective constant and
Theorem~\ref{thm:domain} supplies the common domain. These statements
transfer a favourable return if one is found; they do not assert
unbounded $G_x/\tau_x$. Theorems~\ref{thm:rank}--\ref{thm:barrier}
explain an obstruction to using bounded-order spectral certificates
at subquadratic budgets. Section~\ref{sec:finite} records exact finite
minima, including the full $2+3^{47}$ certificate.

\subsection{Position relative to earlier work}
The derivative lattice, its norm, the additive condition and the
Wronskian are inherited from Pasten \cite{Pasten}; in particular his
Lemma 3.5 concerns nonvanishing, which is weaker than fixing the
positive image generator. The row-index calculation is a specialization
of determinantal divisors and Smith normal form
\cite[Theorem 2.4]{Stanley}. The complementary-weight residue and
composition mechanism is present in Garc\'ia-S\'anchez, Ojeda and
Rosales \cite[Proposition 1(c) and Example 12]{UniqueBetti}, and in
the numerical-semigroup literature
\cite[Section 3.1]{Autry}\cite[Proposition 5.5]{CyrusianKaplan}.
We make that mechanism explicit for signed original-coordinate bounds,
including primes dividing their exponents. The covering proposition
in Appendix~\ref{app:radius} belongs to prescribed-sum rounding;
Cont and Heidari \cite{ContHeidari} provide related optimization results.

The precise claims advanced here are the arithmetic cost formula,
the simultaneous-return comparison with its explicit constants,
the arithmetic growth statement and its stated certificate obstruction,
and the exact original-coordinate examples. Chinese remaindering,
valuation formulas, cyclotomic order arguments and positivity of Fourier
autocorrelations are classical ingredients. The density-domain input
is attributed to Pascadi \cite{Pascadi}, including the use of the
strict margin in his weighted argument; the deduction is written out
in Appendix~\ref{app:domain}. The comparisons above delimit the claims;
they do not constitute a guarantee of bibliographic priority.

\section{Image indices and exact bounded fibres}\label{sec:fibres}
\subsection{The image generator}
For \(n>1\), let
\[
 g_n=\gcd_{p\mid n}\left(\frac{n v_p(n)}p\right),\qquad g_1=0.
\]
\begin{lemma}[Integer-row image calculation]\label{lem:image}
Writing \(h_0=\gcd(g_a,g_b,g_c)\), one has
\begin{equation}\label{eq:image}
 \Dgen=\frac{\gcd(cg_ag_b,\;bg_ag_c,\;ag_bg_c)}{h_0}.
\end{equation}
\end{lemma}
\begin{proof}
By B\'ezout's identity the three derivative blocks have independent
images \(g_a\Z,g_b\Z,g_c\Z\). On their product the two rows become
\((g_a,g_b,-g_c)\) and \((-bg_a,ag_b,0)\).
For independent integer rows \(f,w\), choose a unimodular matrix \(U\)
with \(fU=(\gcd(f),0,\ldots,0)\). The image of \(w\) on the integral
kernel is generated by the remaining entries of \(wU\).
The gcd of the two-row minors, invariant under \(U\), is therefore
\(\gcd(f)\) times that generator. The three minors of the displayed
rows give \eqref{eq:image}. This argument also permits the empty
\(a=1\) block.
\end{proof}

\subsection{A prime block with signed coordinate caps}
Fix \(n=\prod_{i=1}^r p_i^{e_i}>1\) with distinct primes \(p_i\). Define
\begin{equation}\label{eq:block-data}
 \rho_i=\gcd(p_i,e_i),\quad \ell_i=p_i/\rho_i,\quad
 a_i=e_i/\rho_i,\quad
 R=\prod_i\ell_i,\quad G=\gcd_i a_i,\quad f_i=a_i/G.
\end{equation}
A prime dividing its own exponent has \(\ell_i=1\); it contributes no
residual modulus. The nontrivial \(\ell_i\)'s are pairwise coprime and
\begin{equation}\label{eq:block-gcd}
 g_n=\frac nR G,\qquad
 \frac{n e_i/p_i}{g_n}=f_i\frac R{\ell_i}.
\end{equation}
Indeed the numbers \(f_iR/\ell_i\) have gcd one: a common prime outside
\(R\) would divide every \(f_i\), whereas a prime \(p_i\mid R\) does not
divide the \(i\)-th number. This also proves the inverses used next exist.

For integers \(B\ge0,t\), write
\[
 \Phi_{n,B}(t)=\#\{x\in[-B,B]^r\cap\Z^r:D_n(x)=t\}.
\]
\begin{proposition}[Bounded block normal form]\label{prop:normal}
If \(g_n\nmid t\), then \(\Phi_{n,B}(t)=0\). Otherwise put \(s=t/g_n\),
and choose \(0\le r_i<\ell_i\) by
\[
 r_i\equiv s(f_iR/\ell_i)^{-1}\pmod{\ell_i},
 \qquad r_i=0\text{ if }\ell_i=1.
\]
Then
\[
 H=\frac{s-\sum_i f_i(R/\ell_i)r_i}{R}
\]
is integral, and every solution is uniquely of the form
\begin{equation}\label{eq:normal}
 x_{p_i}=r_i+\ell_i y_i,\qquad
 \sum_i f_i y_i=H,\qquad y_i\in\Z.
\end{equation}
The original bounds are exactly
\begin{equation}\label{eq:caps}
 L_i=\left\lceil\frac{-B-r_i}{\ell_i}\right\rceil
 \le y_i\le
 U_i=\left\lfloor\frac{B-r_i}{\ell_i}\right\rfloor.
\end{equation}
If these intervals are nonempty, the count equals
\begin{equation}\label{eq:coefficient}
 \Phi_{n,B}(t)=
 [z^{\,H-\sum_i f_iL_i}]
 \prod_i\left(1+z^{f_i}+\cdots+z^{f_i(U_i-L_i)}\right).
\end{equation}
An empty interval or a negative requested exponent gives count zero.
\end{proposition}
\begin{proof}
Divide the block equation by \(g_n\) and reduce modulo each
\(\ell_i>1\). All terms except the \(i\)-th vanish, forcing its residue.
The numerator defining \(H\) is divisible by every nontrivial modulus
and hence by \(R\), also when \(R=1\).
Substitution gives \eqref{eq:normal}, and each step is reversible.
The positive coordinate steps give \eqref{eq:caps}. Shift \(y_i\) by
\(L_i\) and multiply the finite geometric polynomials to obtain
\eqref{eq:coefficient}.
\end{proof}

The weighted equation in \eqref{eq:normal} is indispensable. For
\(n=15,B=1\), the block \(5x_3+3x_5\) attains precisely
\[
 \{-8,-5,-3,-2,0,2,3,5,8\}.
\]
Its gcd is one and its real range is \([-8,8]\), but it does not attain
\(1\). Divisibility and interval membership do not replace bounded
integer feasibility.

\subsection{Assembly of the three source blocks}
Set \(A_n=\sum_{p\mid n}n v_p(n)/p\), with \(A_1=0\).
For an empty block put \(\Phi_{1,B}(0)=1\), with all other values zero;
also interpret \(\Phi_{n,B}\) at a noninteger argument as zero.
\begin{corollary}\label{cor:global}
For any integer \(w\), the exact number of bounded vectors with
\(F=0,W=w\) is
\begin{equation}\label{eq:global}
 M_w(B)=\sum_{t=-BA_c}^{BA_c}
 \Phi_{a,B}\!\left(\frac{at-w}{c}\right)
 \Phi_{b,B}\!\left(\frac{bt+w}{c}\right)\Phi_{c,B}(t).
\end{equation}
Thus \(\mpr\) is the first budget with \(M_{\Dgen}(B)>0\), and \(m\)
is the first with \(M_w(B)>0\) for some nonzero \(w\).
It suffices to inspect
\(|w|\le B(aA_b+bA_a)\) and multiples of \(\Dgen\).
\end{corollary}
\begin{proof}
For \(D_c=t\), solve \(D_a+D_b=t\) and \(aD_b-bD_a=w\):
\(D_a=(at-w)/c\), \(D_b=(bt+w)/c\). The three disjoint source blocks
are independent, and the values of \(t\) partition the fibre.
The stated bound on \(w\) follows from the triangle inequality.
\end{proof}
These are counts of original vectors, including their signs and
nonprimitive multiples. Formula~\eqref{eq:global} is a reference
description, not a polynomial-time complexity claim.

\section{Proof of the prime-family theorem}\label{sec:main-proof}
For every odd prime $q$, the norm-one vector described after
Theorem~\ref{thm:main} proves $m=1$. Under its hypothesis choose a
simple prime $\ell\mid c$, meaning $v_\ell(c)=1$.
Every coefficient of $D_c$ is divisible by $D=c/R_c$.
The normalized coefficient at $\ell$ is $R_c/\ell$.
A common prime divisor of all normalized coefficients would therefore
be some $p\mid R_c$, $p\ne\ell$; but the coefficient
$(R_c/p)v_p(c)$ is a unit modulo $p$. Hence $g_c=D$.
The equation $F=0$ is equivalent to $z_2+z_q\in D\Z$, with every
such source pair liftable by B\'ezout. Writing $z_q=Dt-z_2$ gives
$W=2Dt-cz_2$, whose image is $D\Z$ since $c/D$ is odd.

Suppose first that $c$ is composite. Then $D\le c/3<q$.
The primitive source equation $2z_q-qz_2=D$ makes $z_2$ odd.
For $z_2\ge1$, $z_q\ge(q+D)/2$; for $z_2\le-1$,
$|z_q|\ge(q-D)/2=:B$, with equality only when $z_2=-1,z_q=-B$.
It remains to lift the target $D_c=(D-c)/2$ inside the same cube.

Write $c=\ell v$ and $r=\sum_{p\mid v}v_p(v)\ge1$; here
$v\ge3$ and $\ell\nmid v$. Put $R=R_c$ and $L=(1-R)/2$.
For each active $p\ne\ell$, choose the centered representative of
$(R v_p(c)/p)z_p\equiv L\pmod p$; set inactive coordinates to zero.
The pivot
\[
 z_\ell=\frac{L-\sum_{p\mid R,\ p\ne\ell}(Rv_p(c)/p)z_p}{R/\ell}
\]
is integral by reduction modulo every prime of $R/\ell$, and the
target row equals $DL$. Since $D\le v$,
$B\ge((\ell-1)v-2)/2\ge v-1$; every nonpivot coordinate is at most
$(v-1)/2$. For the pivot,
\[
 |z_\ell|\le\ell |L|/R+
 \ell\sum_{p\mid R,\ p\ne\ell}v_p(c)|z_p|/p
 <\ell(1+r)/2.
\]
If $\ell\ge5$, $v\ge3^r$ and
$\ell(1+r)\le(\ell-1)3^r-2\le(\ell-1)v-2$.
The first inequality holds at $r=1$; its right-hand increment
$2(\ell-1)3^r$ exceeds the left-hand increment $\ell$.
If $\ell=3$, $v\ge5^r$, and the same induction gives
$3(1+r)\le2\cdot5^r-2\le2v-2$.
Thus the pivot also lies in the cube of radius $B$.
This proves the composite formula and the optimizer description.

If $c$ is prime, $D=1$ and $z_c=z_2+z_q$.
The choice $z_2=-1$ gives the stated vector of norm $(q+1)/2$.
Positive odd $z_2$ forces $z_c\ge(q+3)/2$; $z_2\le-3$ forces
$|z_q|\ge(3q-1)/2>(q+1)/2$. This proves uniqueness, also for $q=3$.

Every powerful integer has a unique representation $a^2b^3$ with
$b$ squarefree: the primes in $b$ are precisely those with odd
exponent. The number at most $T$ is bounded by
$\sqrt T\sum_{b\ge1}b^{-3/2}<3\sqrt T$.
For completeness $\pi(Q)\gg Q/\log Q$ follows elementarily from
\[
 \frac{4^n}{2n+1}\le\binom{2n}{n}\le(2n)^{\pi(2n)}:
\]
Legendre's formula bounds the exponent of $p$ in the binomial
coefficient by $\lfloor\log_p(2n)\rfloor$.
Taking $n=\lfloor Q/2\rfloor$ proves that the relative exceptional
proportion is $O(\log Q/\sqrt Q)$. The inequalities in
\eqref{eq:prime-gap} follow from $D\le c/3$ in the composite branch.
This completes the proof of Theorem~\ref{thm:main}.

\begin{remark}[Necessary boundary distinctions]
At $q=3,c=5$, the primitive minimum is $2$, not the composite expression
$1$. At $q=47,c=49$, $D_c=14z_7$ and $\Dgen=7$.
The putative composite cost $20$ forces $z_2=-1,z_q=-20$ and
$D_c=-21$, which is impossible. The vector $(1,27,2)$ has primitive
norm $27$; the source equation excludes any smaller value except the
already impossible source pair. Thus the nonpowerful hypothesis
cannot simply be dropped.
\end{remark}

\section{A radical-small family and its image}\label{sec:extension}
Write $\hh(n)=n/\rad(n)$ for $n\ge1$, with $\hh(1)=1$.
For an integer $x\ge2$, fix
\begin{equation}\label{eq:family}
 N=N(x)=4\rad_{\rm odd}(x+1)=\lcm(4,\rad(x+1)),\quad
 b=x^N-1,\quad C=Nx^{N-1}.
\end{equation}
Here $\rad_{\rm odd}$ omits the prime $2$. The derivative coordinates
of $x^N$ are the prime coordinates of $x$, and $D_{x^N}=C D_x$.
For the triple $1+b=x^N$, the equations are therefore
\begin{equation}\label{eq:two-row}
 D_b=C D_x,\qquad W=D_b.
\end{equation}
All constants and residues below refer to this particular base $x$.

\begin{proposition}[Radical inequality and exact image]\label{prop:radical-image}
For \eqref{eq:family},
\begin{equation}\label{eq:radical}
 x+1\mid\hh(b),\qquad
 \rad(xb)\le\frac{xb}{x+1}<x^N.
\end{equation}
Writing $g_n=\gcd_{p\mid n}(n v_p(n)/p)$, the integer values of $D_x$
allowed by \eqref{eq:two-row} are precisely $\Delta\Z$, where
\begin{equation}\label{eq:delta}
 \Delta=\frac{g_xg_b}{\gcd(g_b,Cg_x)},\qquad \Dgen=C\Delta.
\end{equation}
If $b$ has a simple prime, define, for $p\mid b$,
\[
 s_p=\begin{cases}
 v_p(x^{\ord_p(x)}-1)-1,&p\text{ odd},\\
 v_2(x-1)+v_2(x+1)-2,&p=2,
 \end{cases}\qquad
 H=\prod_{p\mid b}p^{s_p},\quad I=\prod_{\substack{p\mid b\\p\mid v_p(b)}}p.
\]
Then $s_p\ge0$, $v_p(b)=1+v_p(N)+s_p$, and
\begin{equation}\label{eq:delta-local}
 \Delta=\lcm(g_x,HI),\qquad H\mid\Delta,\qquad\hh(x)\le g_x\le\Delta.
\end{equation}
If also $x$ has a simple prime, $\Delta=g_xHI$.
\end{proposition}
\begin{proof}
For odd $p\mid x+1$, the lifting-the-exponent identity gives
$v_p(b)=v_p(x+1)+v_p(N)$; hence $v_p(\hh(b))\ge v_p(x+1)$
exactly when $p\mid N$. For odd $x$, the identity at $2$ is
$v_2(b)=v_2(x-1)+v_2(x+1)+v_2(N)-1$, and $4\mid N$ suffices.
These prove the divisibility and radical estimate. They also show
that among multiples of $4$, the least exponent giving this
\emph{divisibility} is $N(x)$. This is not a claim of a least exponent
for the radical inequality itself.

The two row images are $g_x\Z,g_b\Z$. Their intersection condition
is $D_x\in g_x\Z$ and $CD_x\in g_b\Z$, yielding \eqref{eq:delta}.
For odd $p\mid b$, $\ord_p(x)\mid N$ and is coprime to $p$;
the valuation formula follows by lifting the exponent. The displayed
formula for $2$ follows from $v_2(N)=2$.
If $b$ has a simple prime, its normalized exponent gcd in
\eqref{eq:block-data} is one, so
$g_b=b/\prod_{p\mid b,\ p\nmid v_p(b)}p$.
At a prime $p\mid b$, $p\nmid x$ and $v_p(C)=v_p(N)$.
Thus the exponent contributed by $g_b/\gcd(g_b,C)$ is
$s_p+\boldsymbol{1}_{p\mid v_p(b)}$.
At primes not dividing $b$ it is zero. Combining with $g_x$ proves
the lcm formula. Every coefficient in the $x$ row is divisible by
$\hh(x)$, and $g_x\mid\Delta$. If $x$ has a simple prime its row gcd
is supported on $x$, which is coprime to $b$, proving the last assertion.
\end{proof}

An explicit large phase is available: if $P>x$ divides $x^2+1$, then
$\ord_P(x)=4$, $P\nmid N$, and $v_P(x^2+1)=1$.
Consequently $P$ is simple in $b$. With $Q=(x^2+1)/P$, reduction of
the identity $(X^N-1)=(X^2+1)U(X)$ and its formal derivative at $x$
gives
\begin{equation}\label{eq:large-phase}
 Q z_P\equiv 2xD_x\pmod P.
\end{equation}
Indeed $N x^{N-1}\equiv2xU(x)\pmod P$ and $b/P=Q U(x)$,
with $U(x)$ a unit. In this differentiation $N$ is held fixed.
The prime-factor theorem of Pascadi \cite[Theorem 1]{Pascadi}
supplies infinitely many such $P>x^{1.3}$. A large modulus alone
does not establish a favourable return or a gap in the minima.

\section{Uniform lifting of all simultaneous phases}\label{sec:comparison}
For a row $D_n=\sum_i A_i z_i$, put $S_n=\sum_i A_i$.
Use the normalized data of \eqref{eq:block-data} and write
$A_i/g_n=f_iR/\ell_i$. An index $j$ with $f_j=1$ is called a
\emph{unit anchor}. A simple prime is always a unit anchor.
For $T\in g_n\Z$, let $r_i(T)$ be the centered representative of
\begin{equation}\label{eq:global-residue}
 r_i(T)\equiv(T/g_n)(A_i/g_n)^{-1}\pmod{\ell_i},
\end{equation}
and set $r_i=0$ if $\ell_i=1$. For even modulus use the representative
in $(-\ell_i/2,\ell_i/2]$; either tie has the same absolute value.
These are residues of the full row, prior to any choice of a
source-zero section or projection.

\begin{lemma}[A drift-balanced lift]\label{lem:lift}
Suppose row $n$ has unit anchor $j$, and put $\alpha=T/S_n$.
If $B\ge|\alpha|$ and $B\ge|r_i(T)|$ for every $i\ne j$, there is
an integer lift of $T$ with
\begin{equation}\label{eq:lift-bound}
 |z_i|\le3B\ (i\ne j),\qquad
 |z_j|\le|\alpha|+\frac{p_j}{e_j} V_n(B),\quad
 V_n(B)=\sum_{i\ne j}\frac{e_i}{p_i}
                  \min\{\ell_i/2,2B\}.
\end{equation}
Conversely any lift has norm at least
$\max(|\alpha|,\max_{i\ne j}|r_i(T)|)$.
\end{lemma}
\begin{proof}
For $i\ne j$, choose a representative in $r_i(T)+\ell_i\Z$
nearest to $\alpha$. Its distance to $\alpha$ is at most $\ell_i/2$
and at most $2B$, since $r_i(T)$ is itself an admissible comparison.
The anchor value $z_j=(T-\sum_{i\ne j}A_i z_i)/A_j$ is integral:
after division by $g_n$, its numerator is divisible by every
nontrivial modulus in $R/\ell_j=A_j/g_n$. Subtract the real solution
with all coordinates $\alpha$ to obtain \eqref{eq:lift-bound}.
The lower bounds follow from the positive row coefficients and
from its congruences. The proof also covers $\ell_j=1$.
\end{proof}

For \eqref{eq:family}, choose a unit anchor in each of the rows $x,b$.
Assume that $b$ has a simple prime. List \emph{all} active nonanchor
coordinates of the two rows; their moduli are denoted $p_i$.
Their residues at targets $\Delta,C\Delta$ are $\kappa_i$.
The list may contain zeros. Define
\begin{align}
 \mu&=\max\{\Delta/S_x,C\Delta/S_b\},&
 \rho(u)&=\max_i|\langle u\kappa_i\rangle_{p_i}|,\nonumber\\
 h_x(u)&=\max\{\mu u,\rho(u)\},&
 G_x&=h_x(1),\qquad\tau_x=\min_{u\ge1}h_x(u),\label{eq:return}
\end{align}
where $\langle\cdot\rangle_p$ denotes the centered representative,
and the maximum of an empty phase list is zero.
Let $m_x(u)$ be the least original norm with $D_x=u\Delta$ and
$D_b=uC\Delta$. These are positive image indices; sign symmetry
gives $m(1,b)=\min_{u\ge1}m_x(u)$ and $\mpr(1,b)=m_x(1)$.
The minimum defining $\tau_x$ exists because $\mu>0$.

\begin{theorem}[Uniform original-coordinate comparison]\label{thm:comparison}
Put
\[
 H_0=\hh(x-1)\hh(x+1)\hh(x^2+1),\quad
 \sigma_n=\sum_{p\mid n}1/p,\quad
 \vartheta_n=p_{j_n}/e_{j_n},\quad J=H_0(\sigma_b+1)/2,
\]
where $j_x,j_b$ are the chosen unit anchors, and set
\begin{equation}\label{eq:Cx}
 C_x=\max\{3,\ 1+\vartheta_x(2\sigma_x+J),\
                    1+\vartheta_b(4\sigma_b+2+J)\}.
\end{equation}
For every positive integer $u$,
\begin{equation}\label{eq:pointwise}
 h_x(u)\le m_x(u)\le C_x h_x(u).
\end{equation}
In particular
\begin{equation}\label{eq:compare}
 \tau_x\le m(1,b)\le C_x\tau_x,\quad
 G_x\le\mpr(1,b)\le C_xG_x,\quad
 \frac{G_x}{C_x\tau_x}\le\frac{\mpr(1,b)}{m(1,b)}
 \le\frac{C_xG_x}{\tau_x}.
\end{equation}
If $H_0\le K$, $\sigma_x+\sigma_b\le L$ and
$\vartheta_x,\vartheta_b\le Y$, one may use the base-independent
constant
\begin{equation}\label{eq:C0}
 C_0=\max\{3,\ 1+Y(4L+2+K(L+1)/2)\}.
\end{equation}
No bound on $\Delta$ or on the sum of moduli is assumed.
\end{theorem}
\begin{proof}
Write $t=\log_2\Delta$ and $\lambda_b=\sum_{p\mid b}v_p(b)/p$.
Equation~\eqref{eq:delta-local}, and $v_p(N)=1$ at odd primes
dividing $N$ while $v_2(N)=2$, imply
\begin{equation}\label{eq:lambda-est}
 \lambda_b\le2\sigma_b+1+t/2.
\end{equation}
Indeed $\sum s_p\le t$ and $p\ge2$.
Since $N\hh(x+1)\ge2(x+1)$, also $NH_0\ge2(x+1)$.
As $b<x^N$ and $S_b=b\lambda_b$,
\[
 \mu\ge\frac{C\Delta}{S_b}>\frac{2\Delta}{H_0\lambda_b}.
\]
For $t\ge0$, $t/2^t\le1$ and $t^2/2^t\le2$; these follow by
maximizing the two elementary real functions. Consequently
\[
 \frac{t}{2\mu}<\frac{H_0t(2\sigma_b+1+t/2)}{4\,2^t}
 \le H_0(\sigma_b+1)/2=J
\]
when $t>0$, and the same non-strict bound holds at $t=0$.

For the source row, split each $e_p$ into $1+(e_p-1)$.
The first part contributes at most $2\sigma_xB$ to $V_x(B)$,
and the second at most $\tfrac12\sum(e_p-1)\le t/2$, since
$\hh(x)\le\Delta$. For the target row use
$e_p=1+v_p(N)+s_p$: the first two parts contribute at most
$(4\sigma_b+2)B$, including the prime $2$, and the last at most $t/2$.
Thus, whenever $B\ge\mu$,
\[
 V_x(B)\le(2\sigma_x+J)B,\qquad
 V_b(B)\le(4\sigma_b+2+J)B.
\]
Apply Lemma~\ref{lem:lift} to both rows at $B=h_x(u)\ge\mu$.
Their coordinates are disjoint, and their targets use the same $u$.
This proves \eqref{eq:pointwise}. Taking minima and dividing the
appropriate upper and lower bounds gives \eqref{eq:compare}.
The constant occurs once, not squared. Formula~\eqref{eq:C0} follows
by enlarging the three parameters.
\end{proof}

\begin{theorem}[A common positive-density domain]\label{thm:domain}
Fix $1<\alpha<1.3$. There are finite constants $K,L,Y$ such that
a set of bases of positive lower density satisfies all of the following:
\begin{enumerate}
 \item $H_0\le K$ and $\sigma_x+\sigma_b\le L$;
 \item $x$ and $b=x^{N(x)}-1$ each have a simple prime at most $Y$;
 \item $x^2+1$ has a prime divisor $P>x^\alpha$.
\end{enumerate}
Hence \eqref{eq:compare} holds there with the same $C_0$.
\end{theorem}
The proof, with the analytic input and the order of parameter choices
specified, is given in Appendix~\ref{app:domain}. It uses subtraction
of exceptional sets of arbitrarily small upper density from a
positive-density prime-factor set. It does not intersect unrelated
positive-density sets without further information.

\subsection{What an unbounded separation would require}
On any domain with fixed $C_0$, the ratio $\mpr/m$ is unbounded along
a sequence if and only if $G_x/\tau_x$ is. A sufficient and necessary
return formulation along such a sequence is the existence of indices
$u=u(x)$ with
\begin{equation}\label{eq:required-return}
 \mu_x u/G_x\longrightarrow0,\qquad
 \max_i|\langle u\kappa_i(x)\rangle_{p_i(x)}|/G_x\longrightarrow0.
\end{equation}
All phases in this display belong to the same base and the same index.

\begin{proposition}[First-wrap obstruction]\label{prop:wrap}
For a nonzero centered phase $\theta=|\kappa|$ modulo $p$,
\[
 \tau_x\ge\min\{\theta,\mu p/(\theta+\mu)\}.
\]
If $G_x=\mu$, then $\tau_x=G_x$ and $\mpr/m\le C_x$.
Otherwise let $p_*$ be the largest modulus among phases with
$|\kappa_i|=G_x$. Then
\begin{equation}\label{eq:wrap-capacity}
 \frac{\mpr}{m}\le C_x\max\{1,G_x(G_x+\mu)/(\mu p_*)\}.
\end{equation}
\end{proposition}
\begin{proof}
If $h_x(u)=B<\theta$, the residue bound cannot hold before the first
wrap: $u\theta\ge p-B$. Since $u\le B/\mu$, this forces
$B\ge\mu p/(\theta+\mu)$. Otherwise $B\ge\theta$ already.
Apply this to $\theta=G_x$, or use $h_x(u)\ge\mu u$ when $G_x=\mu$,
and then \eqref{eq:compare}.
\end{proof}
In addition $\tau_x\ge\mu$, so unbounded separation with fixed $C_0$
requires both $G_x/\mu$ and the second capacity in
\eqref{eq:wrap-capacity} to become unbounded. These conditions are
necessary, not sufficient.

\section{Growing arithmetic dimension and sparse certificates}\label{sec:growth}
This section concerns the family \eqref{eq:family}. Density statements
use ordinary counting of integer bases: an exceptional set has size
$o(X)$ up to $X$, even when the ambient domain itself is restricted.
They do not assert independence of its modular phases.

\begin{lemma}[Cyclotomic blocks]\label{lem:cyclo}
Let $x\ge5$, $n=\rad_{\rm odd}(x+1)$, and $d\mid n$, $d\ge3$.
If $(2d+1)^2>x+1$, each of the three orders $d,2d,4d$ is realized
by a distinct active prime of $b=x^{4n}-1$. These primes are at least
$2d+1,2d+1,4d+1$, respectively. An active prime means
$p\nmid v_p(b)$.
\end{lemma}
\begin{proof}
For an odd prime $p\nmid x$, put $f=\ord_p(x)$ and
$a=v_p(x^f-1)$. The factorization into cyclotomic polynomials and
the valuation formula $v_p(x^{ft}-1)=a+v_p(t)$ give
\[
 v_p(\Phi_m(x))=
 \begin{cases}a,&m=f,\\1,&m=fp^j\ (j\ge1),\\0,&\text{otherwise}.
 \end{cases}
\]
This follows by induction on the divisors in
$x^m-1=\prod_{e\mid m}\Phi_e(x)$.
For $p=2$, an index $m\ge3$ can contribute only when $m$ is a power
of two, in which case the valuation is one. For the three indices
here this does not occur. If $m=fp^j$, then $f\mid p-1$, so $p$
is the largest prime factor of $m$. Thus the product of the
nonprimitive prime contributions to $\Phi_m(x)$ is at most $m$.
But the root factorization gives
\[
 (x-1)^{\varphi(m)}\le\Phi_m(x)\le(x+1)^{\varphi(m)}.
\]
For $m=d,2d$, $\varphi(m)\ge2$ and $(x-1)^2>2(x+1)$;
for $m=4d$, $\varphi(m)\ge4$ and $(x-1)^4>4(x+1)$.
Since $d\le x+1$, these bounds exceed $m$ and force a primitive
prime with exact order $m$.

Such a prime satisfies $p\equiv1\pmod m$; parity improves the
lower bound at odd order $d$ to $2d+1$. It does not divide $N$:
odd primes dividing $N$ divide $x+1$ and have order at most two.
Its exponent in $b$ is therefore its exponent in $\Phi_m(x)$.
In all three cases $\varphi(m)<p/2$, while $p^2>x+1$; hence
$(x+1)^{\varphi(m)}<p^p$. Its positive exponent is less than $p$,
so it is active. Exact orders make all primes from distinct indices
different; odd $d$ also separates the three types of index.
\end{proof}

\begin{theorem}[Growing active rank in density]\label{thm:rank}
Fix $K\ge1$, $0<\beta<1$ and an integer $q\ge1$. Outside $o(X)$
exceptions among $x\le X$ with $H_0(x)\le K$, the row of $b$
has at least $q$ distinct active primes greater than $x^\beta$.
After deleting any fixed number of coordinates, the product of
remaining active moduli exceeds every fixed power of $x$ in the
same density sense.
\end{theorem}
\begin{proof}
For fixed $Y$, let $Z_Y(x)$ count the odd primes $p\le Y$ dividing
$x+1$, and let $A_Y=\sum_{3\le p\le Y}1/p$.
CRT over a fixed finite set gives mean $A_Y$ and variance at most
$A_Y$ in the limiting density. Since $A_Y\to\infty$, for each fixed
$s$ the upper density of $Z_Y<s$ is at most
$A_Y/(A_Y-s)^2$, tending to zero with $Y$.
If $Z_Y\ge s$, choose $s$ of those primes and let their product be
$Q_0\mid Q_Y=\prod_{3\le p\le Y}p$. Since
$n\ge(x+1)/(2K)$, every $t\mid Q_0$ gives a distinct divisor
$d=n/t\ge(x+1)/(2KQ_Y)$.
With $K,Y$ fixed, for sufficiently large $x$ these $2^s$ divisors
satisfy the previous lemma and give $3\cdot2^s$ active primes
greater than $x^\beta$. First choose $s$ for the requested $q$,
then let $Y$ grow to make the upper exceptional density arbitrarily
small. This proves a zero-density exception. For the product claim
choose $q$ larger than the number removed plus the requested power
divided by $\beta$. No growing-$q$ rate is asserted.
\end{proof}

\begin{lemma}[Simple phases at feasible short budgets]\label{lem:short}
Fix $K,Y$, $0\le\theta<\beta<2$ and $q\ge1$. Outside $o(X)$
exceptions among bases $x\le X$ with $H_0\le K$ and simple anchors
at most $Y$ in both rows, the following alternative holds:
either there are no real budgets $B$ and integer horizons $H\ge1$
with $\mu H\le B\le x^\theta$, or the global phase list contains at
least $q$ distinct nonzero simple nonanchor phases with modulus
$p>x^\beta$. Each has $p\nmid\Delta$ and $p>2B$ for all large $x$.
\end{lemma}
\begin{proof}
We give the size estimates to specify the quantifiers. First
\begin{equation}\label{eq:logdrift}
 \lambda_b<4\log_2(2x\Delta),\qquad
 \mu>\frac{\Delta}{2K\log_2(2x\Delta)}
       \ge\frac1{2K\log_2(2x)}.
\end{equation}
Indeed $\hh(b)=\prod p^{v_p(N)+s_p}$ divides $N H$, so is at most
$N\Delta$. Split $\lambda_b=\sigma_b+\sum(e_p-1)/p$.
With $A=\log_2 b$, primes $p\le A$ contribute at most
$1+\log A$ to $\sigma_b$, and those above $A$ contribute at most
$\omega(b)/A\le1$; also
$\sum(e_p-1)/p\le\tfrac12\log_2(N\Delta)$.
Using $N\le4(x+1)$ gives the first (loose) bound for $x\ge5$.
The earlier lower bound $\mu>2\Delta/(K\lambda_b)$ proves the
second. The last inequality follows from
$\log_2(2x\Delta)\le\Delta\log_2(2x)$ for $\Delta\ge1$.
Thus feasibility forces $\Delta<x^\gamma$ eventually for every
fixed $\gamma>\theta$: otherwise the increasing function
$t/\log_2(2xt)$ at $t=x^\gamma$ would contradict $\mu\le x^\theta$.

We strengthen the cyclotomic size argument on the domain
$\Delta\le x^D$, for any fixed $D$. Choose the large orders
$m=d,2d,4d$ in the proof of Theorem~\ref{thm:rank}, so $m\gg x$
with constants fixed before $x$ grows. Suppose all primitive primes
of $\Phi_m(x)$ are at most $Z=x^\beta$.
There are at most $\lfloor(Z-1)/m\rfloor$ such primes, since each
is $1\pmod m$. Their repeated parts divide $H\mid\Delta$,
and the nonprimitive part is at most $m$. Hence
\begin{equation}\label{eq:block-size}
 (x-1)^{\varphi(m)}\le\Phi_m(x)
       \le m\Delta Z^{\lfloor(Z-1)/m\rfloor}.
\end{equation}
For every $\varepsilon>0$ one has
$\varphi(m)\ge c_\varepsilon m^{1-\varepsilon}$: in the product
for $\varphi(m)m^{\varepsilon-1}$, only finitely many primes can
give a factor smaller than one. Choose $\varepsilon<2-\beta$.
The logarithm of the left side of \eqref{eq:block-size} is
$\gg x^{1-\varepsilon}\log x$, whereas that of the right side is
$O(\log x+x^{\beta-1}\log x)$, a contradiction.
Thus every one of arbitrarily many fixed large orders supplies a
prime $p>x^\beta$. At most $\lfloor D/\beta\rfloor$ of these
distinct primes can be repeated, because then each divides
$H\mid\Delta\le x^D$. Start with enough orders to leave $q$ simple
ones. The fixed-$Y$ CRT argument again makes the exceptional count
$o(X)$, without asserting any density for the feasible branch.

Choose $\theta<\gamma<\beta$ and use feasibility to put
$\Delta<x^\gamma$. At the remaining simple primes, $p\nmid\Delta$
and eventually $p>Y,2B$. To check that their \emph{global} phase is
nonzero, note that with simple primes in both rows,
\[
 \Delta=\epsilon_0 g_xg_b/N,\qquad
 C\Delta/g_b=\epsilon_0g_xx^{N-1},\qquad
 \epsilon_0=\begin{cases}1,&x\text{ odd},\\4,&x\text{ even}.
 \end{cases}
\]
Here $g_x$ is supported on $x$, and
$\gcd(g_b,C)=N/\epsilon_0$ by the local valuations in
Proposition~\ref{prop:radical-image}. The right side is a unit at
each new prime $p$, as is its normalized row coefficient. Its
residue \eqref{eq:global-residue} is therefore nonzero.
\end{proof}

\subsection{What positivity can certify}
Let a finite list of phases be $(p_i,\kappa_i)$, let $H\ge1$,
$M=H+1$, and choose integer lengths $1\le L_i\le p_i$.
On $\Z/p_i\Z$ set
\[
 f_i(z)=\frac1{L_i}\#\{(a,b)\in\{0,\ldots,L_i-1\}^2:a-b=z\},
 \quad V_L(t)=\left|\frac1L\sum_{a=0}^{L-1}e^{2\pi iat}\right|^2,
 \quad D_i=L_i/p_i.
\]
Thus $0\le f_i\le1$ and $f_i(0)=1$.

\begin{proposition}[Positive Fourier identity and finite caps]\label{prop:spectrum}
The energy
\begin{align}
 \mathcal E&=\frac1M\sum_{a,b=0}^H\prod_i f_i((a-b)\kappa_i)\nonumber\\
 &=M\prod_iD_i\sum_{m_i\bmod p_i}
 V_M\!\left(\sum_i\frac{m_i\kappa_i}{p_i}\right)
             \prod_iV_{L_i}(m_i/p_i)\label{eq:energy}
\end{align}
has nonnegative summands and diagonal contribution one in its first
expression. Therefore a lower certificate $\mathcal E>1$ yields
some $1\le u\le H$ with
$|\langle u\kappa_i\rangle_{p_i}|\le L_i-1$ for every $i$.
For any $q$ designated coordinates with $D_i\le\delta$, the total
mass of all frequency terms involving at most $r\le q$ of these
coordinates is at most
\begin{equation}\label{eq:support-cap}
 M\binom qr\delta^{q-r}.
\end{equation}
Any $T$ individually listed distinct frequency terms have mass at most
\begin{equation}\label{eq:term-cap}
 TM\delta^q.
\end{equation}
These bounds allow arbitrary frequencies in the undesignated coordinates.
\end{proposition}
\begin{proof}
Finite Fourier inversion gives $\widehat f_i(m)=D_iV_{L_i}(m/p_i)$.
The normalized sum over $a,b$ gives $M V_M$ and proves the identity.
If the energy exceeds its diagonal, an off-diagonal product is positive,
so each difference is represented by two interval elements, giving
the claimed centered bound.

Fix a subset $J$ of $r$ designated coordinates and set the other
$q-r$ frequencies to zero. Removing their factors $D_i$ leaves the
full Fourier energy of the remaining coordinates, at most $M$ by
the time expression and $f_i\le1$. Cover all supports of size at
most $r$ by the $\binom qr$ choices of $J$ and use positivity;
this proves \eqref{eq:support-cap}, even though the cover can overlap.
For a single term $V_L\le1$ and each undesignated $D_i\le1$,
so it is at most $M\delta^q$. Summing $T$ terms proves the second cap.
\end{proof}

\begin{theorem}[A restricted spectral-certificate barrier]\label{thm:barrier}
Fix $K,Y$, $0\le\theta<2$, an interaction order $r\ge0$ and
$c\ge0$. Consider bases with $H_0\le K$ and simple anchors at most
$Y$, and feasible budgets $\mu H\le B\le x^\theta$, $H\ge1$.
Use all global phases and $L_i=\min(p_i,\lfloor B\rfloor+1)$.
Outside $o(X)$ bases up to $X$, uniformly over these feasible choices,
neither the complete collection of terms of fixed interaction order
$r$ nor any at most $x^c$ individually listed distinct terms has mass
exceeding one in \eqref{eq:energy}; in fact their masses tend to zero.
For the first assertion it suffices to count interactions within an
appropriate fixed number of designated large coordinates.
\end{theorem}
\begin{proof}
If $\theta=0$, \eqref{eq:logdrift} still gives the following bound;
otherwise the same calculation applies. Fix $\theta<\beta<2$.
By Lemma~\ref{lem:short} choose any required fixed number $q$ of
nonzero simple phases $p_i>x^\beta$, outside $o(X)$ exceptions.
Then $D_i\le2x^{\theta-\beta}$ and
$M=O_K(x^\theta\log x)$ by \eqref{eq:logdrift}.
Choose $q>r+\theta/(\beta-\theta)$ in \eqref{eq:support-cap},
and $q>(c+\theta)/(\beta-\theta)$ in \eqref{eq:term-cap}.
Both bounds tend to zero. The number $q$ is fixed before the density
limit, and the bounds use only the displayed budget inequalities,
so are uniform in $B,H$.
\end{proof}
This is a limit on two specified positive-sum certificates. It is
not an upper bound for the full energy, an absence-of-returns theorem,
or a general algorithmic lower bound. It leaves open exceptional bases,
budgets $o(G_x)$ larger than a fixed subquadratic scale, and collective
frequency mass with increasing support. Lemma~\ref{lem:short} does
not assert that feasible subquadratic budgets have positive density.

\section{Finite certificates in the original coordinates}\label{sec:finite}
The ordinary proofs above establish their universal statements.
The finite results in this section additionally use supplied exact data
and a short standard-library verifier. No floating-point optimization,
probable-prime declaration or assumed factorization is used in their
acceptance checks.

\begin{theorem}[Two variable-base minima]\label{thm:finite}
For $b=x^{60}-1$, with the original coordinates and primitive meaning
of \eqref{eq:minima}, the following minima hold:
\[
\begin{array}{c|r|r|r|r}
 x&\Delta&u\text{ at a useful minimizer}&m(1,b)&\mpr(1,b)\\\hline
14&1&331476352&1054020192&311136274396595031\\
44&36&1807554653&70930543252&92514302310809204735
\end{array}
\]
Here $N(x)=60$ in both cases, so both triples satisfy
$\rad(xb)<x^{60}$. The ratios are approximately
$2.95\cdot10^8$ and $1.30\cdot10^9$, respectively.
\end{theorem}
\begin{proof}[Certificate verification]
The file \path{companion/data/exact_minima.json} gives complete
factorizations, row anchors, integer vectors and lattice bases.
Its prime certificates are in \path{companion/data/minima_primes.json}.
The program
\texttt{companion/verify\_finite.py} performs the following exact
certificate proof, whose exhaustiveness we specify here.

Each prime $p>2$ is certified by a complete factorization of $p-1$
into previously certified primes and an integer $a$ such that
$a^{p-1}=1\pmod p$ and
$\gcd(a^{(p-1)/r}-1,p)=1$ for every prime $r\mid p-1$.
For any prime divisor $s\mid p$, these conditions force the order
of $a\pmod s$ to be $p-1$, so $s\ge p$ and $p$ is prime.
Multiplication then verifies each whole factorization and reconstructs
all original rows, their image and the global phases.

For a proposed lower budget $B=m-1$, a possible positive index obeys
$1\le u\le\lfloor B/\mu\rfloor$ and every centered phase is at most
$B$. Choose the specified pivot $(p,\kappa)$ with $p>2B$.
A displayed integer basis $(q,r),(s,t)$ has determinant $p$ and lies
in the congruence lattice $z=\kappa u\pmod p$; it is therefore a
basis of the whole lattice. Write
$(u,z)=a(q,r)+v(s,t)$. The inverse coordinate
$v=(qz-ru)/p$ has its extrema at the four rectangle corners
$1\le u\le\lfloor B/\mu\rfloor$, $-B\le z\le B$.
Enumerating every integer $v$ in this interval and the corresponding
complete interval of $a$ lists every possible index. Every remaining
global phase is then tested, not just the pivot phase.

For $x=44$ no index survives. For $x=14$ the sole survivor is
$u=331476352$. Proposition~\ref{prop:normal}, applied to both
original cubes, rejects it: the target-row carry is $1420612936$,
whereas the upper possible carry is $1420612934$, a deficit of two.
In these two inputs the normalized weights are only one or two.
Their interval sum is the whole interval from its lower to upper
endpoint if any weight-one coordinate is flexible; otherwise it is
the corresponding parity class. This checks the carries exhaustively.
The displayed vectors attain $m$ at the asserted indices.

At $u=1$, the computed lower bound $\lceil G_x\rceil$ equals the
listed $\mpr$ in each case, and the supplied primitive vectors attain
it. All row equations, coordinates, norms and image values are checked
as integers. These prove both minima, not only upper bounds.
\end{proof}

For $x=14$, the winning index is $4\pmod7$. A section forcing the
nonanchor source coordinate to zero would require $u=0\pmod7$ and
miss it. The return minimum also differs from the original minimum:
$\lceil\tau_{14}\rceil=m-5$, while $\tau_{44}=m$.
These are concrete reasons to retain the source phases and the final
integer carry when using a reduced return problem.

\subsection{Finite mixed-frequency certificates}
Two companion benchmarks illustrate the role of mixed frequencies
in Proposition~\ref{prop:spectrum}. The variable-base case
$x=6,N=28$ at return budget $B=11000,H=1307$ has
\[
 \mathcal E=3440888017/2931421802>1,
 \qquad\mathcal E_{\rm axes}=41388640220693/61401273111326<1.
\]
The fixed-power case $k=53$, on its specified normalized return
section at $B=2000000,H=400000$, has
\[
 \begin{aligned}
 \mathcal E&=\frac{2073236796034760609}{1600005600004400001}>1,\\
 \mathcal E_{\rm axes}
   &=\frac{91767038289310419629485}{156111755376230439915667}<1.
 \end{aligned}
\]
The first uses 23 explicit frequency terms, 20 mixed; the second
uses 11, 10 mixed. Their certified lower scores exceed $1281/1250$
and $5249/5000$. The axis mass means all terms supported on either
of the two restrictive axes, with the zero term counted once.
All other phases are automatically paid by the declared budget.
The data include original-coordinate witnesses and their separate
norm bounds; a return budget is not identified with an original norm.

For reproducibility, the verifier both sums the rational time-domain
energy and checks rational lower bounds for the selected frequencies.
If $a=\operatorname{dist}(t,\Z)$, $d=\operatorname{dist}(Lt,\Z)$
and $a\ne0$, a lower bound for $V_L(t)$ is the maximum of zero,
\[
 1-\frac{10}{3}(L^2-1)a^2,\qquad
 \left[\frac{d}{La}
 \left(1-\frac53d^2+\frac56d^4-\frac{25}{126}d^6\right)\right]^2.
\]
At $a=0$ the value is one. The first estimate follows from the
cosine lower bound in the squared exponential sum and $\pi^2<10$.
For the second use $\sin(\pi a)\le\pi a$ and the alternating
seventh-degree sine lower bound on $0\le\pi d\le\pi/2$;
its bracket decreases with $\pi^2$ on this range, so substitution
of $10$ gives the displayed smaller positive polynomial.
These finite successes concern two restrictive phases. They do not
contradict the growing-dimensional statement of Theorem~\ref{thm:barrier}.

\section{A compact exact certificate for \texorpdfstring{\(2+3^{47}\)}{2 plus 3 to the power 47}}\label{sec:k47}
The large example can be checked without enumerating a large cube.
Its certificate consists of a factorization, a two-dimensional lattice
basis, a small affine deficit and a primitive congruence.

\begin{theorem}\label{thm:k47}
For \(a=2,b=3^{47}\), in the norm of \eqref{eq:minima},
\[
 m=6493223,\qquad \mpr=1358381707976.
\]
There are exactly twenty vectors attaining \(m\), with Wronskian
values \(\pm133972\Dgen\).
All vectors attaining \(\mpr\) with \(W=\Dgen\) are parametrized in
\eqref{eq:primitive-family} below.
This is a radical-small triple:
\(\rad(abc)/c=6/343\).
\end{theorem}

\subsection{Model and integral projected coordinates}
Use the factorization
\begin{equation}\label{eq:k47-factor}
 3^{47}+2=7^4\cdot11\cdot59\cdot4691\cdot3637447066471.
\end{equation}
Each displayed factor is prime; the companion verifies this by exhaustive
trial division up to its integer square root, and verifies the product.
Put
\[
 \begin{gathered}
 P=3637447066471,\quad d=7\cdot11\cdot59\cdot4691=21311213,\quad h=343,\\
 T=3^{46}=8862938119652501095929,\qquad c=hdP=3T+2,\\
 L=d+2\left(4d/7+d/11+d/59+d/4691\right)=50273151.
 \end{gathered}
\]
The coordinate order is \((2,3,7,11,59,4691,P)\).
The rows are \(F=x_2+47Tx_3-D_c\) and
\(W=T(94x_3-3x_2)\).
Since \(g_c=h\), formula~\eqref{eq:image} gives \(\Dgen=Th\).

For any integral additive vector set
\[
 j=(3x_2-94x_3)/h,\qquad v=x_P.
\]
These are integral: reducing \(3F=0\) modulo \(h\), and using
\(3T\equiv-2\pmod h\), proves the divisibility in \(j\).
Elimination gives \(W=-Thj\) and
\begin{equation}\label{eq:projected}
 j\equiv3dv\pmod P,\qquad
 |v|\le B,\qquad |Tj+2dv|\le PLB
\end{equation}
whenever \(\norm{x}\le B\).
Indeed,
\(Tj/(dP)=x_2-2(4x_7/7+x_{11}/11+x_{59}/59+x_{4691}/4691)-2v/P\).
Only necessity of these real inequalities is used in the exclusion.

\subsection{Only one useful projected pair at the proposed budget}
Let \(B=6493223\) and, by negation, assume \(j>0\).
The congruence lattice
\(\Gamma=\{(j,v)\in\Z^2:j\equiv3dv\pmod P\}\) has index \(P\).
The columns
\[
 u=(133972,1649929),\qquad z=(-2158411,568941)
\]
belong to \(\Gamma\) and have determinant \(P\), hence form a basis.
Their values of \((3dv-j)/P\) are \(29\) and \(10\).

Equation~\eqref{eq:projected} implies
\[
 1\le j\le\left\lfloor\frac{(PL+2d)B}{T}\right\rfloor=133972.
\]
If \((j,v)=\alpha u+\beta z\), then
\[
 \beta=\frac{133972v-1649929j}{P},\qquad
 |\beta|\le\frac{133972(B+1649929)}P<1.
\]
Therefore \(\beta=0\) and \(\alpha=1\).
Every useful vector of norm at most \(B\), up to negation, projects to
\(u\). The same cover applies to every smaller budget.

\subsection{The affine phase and all twenty minimizers}
Write \((p_i)=(7,11,59,4691)\) and \((e_i)=(4,1,1,1)\).
For a projected pair define
\[
 t=\frac{3dv-j}{P},\qquad
 r_i\equiv-t\bigl(3e_i(d/p_i)\bigr)^{-1}\pmod{p_i},
 \quad 0\le r_i<p_i.
\]
Substitution in \(F=0\) reduces it to
\[
 47x_3=H_0+3\sum_i e_i y_i,\qquad
 H_0=\frac{t+\sum_i3e_i(d/p_i)r_i}{d},\qquad x_{p_i}=r_i+p_i y_i.
\]
The residues make \(H_0\) integral.
Choose \(0\le x_0<3\) with \(47x_0\equiv H_0\pmod3\), and put
\(\theta=(47x_0-H_0)/3\),
\(C_0=(hj+94x_0)/3\).
This last number is integral: modulo three, \(dP\equiv-1\) and
\(Pt\equiv-j\) give \(H_0\equiv j\), while
\(47x_0\equiv H_0\) gives \(x_0\equiv2j\).
Every lift is uniquely
\begin{equation}\label{eq:k47-affine}
 x_2=C_0+94n,\quad x_3=x_0+3n,\quad
 x_{p_i}=r_i+p_i y_i,\quad x_P=v,\quad
 4y_7+y_{11}+y_{59}+y_{4691}=\theta+47n.
\end{equation}
The derivation reverses, proving sufficiency as well.

For \(u\), the data are
\[
 t=29,\quad (r_i)=(5,2,47,2335),\quad H_0=13,\quad
 x_0=2,\quad \theta=27,\quad C_0=15317528.
\]
Let
\[
 \ell_i(B)=\left\lceil\frac{-B-r_i}{p_i}\right\rceil,\quad
 U_i(B)=\left\lfloor\frac{B-r_i}{p_i}\right\rfloor,\quad
 \mathcal L(B)=\sum_i e_i\ell_i(B),\quad
 N_+(B)=\left\lfloor\frac{B-C_0}{94}\right\rfloor.
\]
Every bounded lift must have
\(\mathcal L(B)\le27+47n\) and \(n\le N_+(B)\).
The exact first-contact calculation is
\[
\begin{array}{r|r|r|r}
 \text{budget}&\mathcal L&N_+&27+47N_+\\ \hline
 6493222&-4412144&-93876&-4412145\\
 6493223&-4412148&-93876&-4412145
\end{array}
\]
The first row excludes every lift at the preceding budget.
At \(B=6493223\) the lower bounds are
\[
 (\ell_i)=(-927604,-590293,-110055,-1384).
\]
They and the source bound force \(n=-93876\).
Putting \(y_i=\ell_i+d_i\) leaves
\[
 4d_7+d_{11}+d_{59}+d_{4691}=3,\qquad d_i\ge0.
\]
All upper caps exceed three. Hence \(d_7=0\), and the remaining
three deficits are all nonnegative compositions of three.
There are \(\binom52=10\). The resulting vectors are
\begin{equation}\label{eq:20}
 \begin{aligned}
 x_2&=6493184,&x_3&=-281626,&x_7&=-6493223,\\
 x_{11}&=-6493221+11d_{11},&
 x_{59}&=-6493198+59d_{59},\\
 x_{4691}&=-6490009+4691d_{4691},&
 x_P&=1649929,
 \end{aligned}
\end{equation}
where \(d_{11}+d_{59}+d_{4691}=3\).
Each satisfies the original rows and has norm \(B\).
Their negatives give the other ten; the two lists are disjoint because
\(W\ne0\). The complete projected cover and preceding-budget exclusion
prove the useful minimum and completeness of this list.

\subsection{The primitive minimum and its entire optimizer fibre}
First use \(j=1\), so \(W=-\Dgen\).
The congruence \(3dv\equiv1\pmod P\) has centered representative
\[
 v_0=-1358381707976,\qquad 2|v_0|<P.
\]
Every such vector has norm at least \(M=|v_0|\); within budget \(M\)
its pivot is necessarily \(v_0\).
The negative primitive fibre has
\[
 t=-23875615,\quad(r_i)=(2,8,25,1937),\quad
 H_0=7,\quad x_0=2,\quad\theta=29,\quad C_0=177.
\]
In particular
\[
 (-177,-2,-2,-327,-25,-1937,1358381707976)
\]
has \(F=0\), \(W=\Dgen\) and norm \(M\), proving the primitive minimum.

For completeness, all positive primitive minimizers are precisely the
negatives of
\begin{equation}\label{eq:primitive-family}
 \begin{gathered}
 (177+94n,\ 2+3n,\ 2+7y_7,\ 8+11y_{11},\
     25+59y_{59},\ 1937+4691y_{4691},\ -M),\\
 4y_7+y_{11}+y_{59}+y_{4691}=29+47n,\qquad
 n,y_7,y_{11},y_{59},y_{4691}\in\Z,
 \end{gathered}
\end{equation}
subject to absolute value at most \(M\) in each coordinate.
This is a bounded affine parametrization, not an enumeration.
Its completeness follows from the unique pivot representative and
\eqref{eq:k47-affine}. The last coordinate ensures every admitted vector
has norm exactly \(M\).
Finally \eqref{eq:k47-factor} gives \(\rad(abc)/c=6/343\).
The ratio of the two costs exceeds \(200000\), but a single example,
however large, proves no unboundedness in a restricted infinite family.

\section{Open questions, scope and reproducibility}\label{sec:limits}
\begin{question}\label{q:power}
Is $\mpr(2,3^k)/m(2,3^k)$ unbounded as $k$ ranges over positive integers?
\end{question}
\begin{question}\label{q:radical}
Is $\mpr(a,b)/m(a,b)$ unbounded over coprime triples $a+b=c$ with
$\rad(abc)<c$?
\end{question}
Neither question is settled here. Theorem~\ref{thm:main} supplies
unbounded separation in a different family. The fixed-power question
has the necessary simple-prime and non-Wieferich consequence in
Theorem~\ref{thm:fixed}. Conditional non-Wieferich results, such as
\cite[Theorems 1.2--1.4 in the cited arXiv version]{FelliniMurty},
do not give that consequence unconditionally or supply the returns.
For the variable-base route to Question~\ref{q:radical},
\eqref{eq:required-return} is the precise remaining existence target.
The common domain and the comparison theorem remove a lifting
obstruction but do not provide this sequence.

The arithmetic growth and Fourier bounds identify limitations of a
specific certificate format. Larger finite records alone do not prove
an infinite gap; an independent supply of simultaneous returns is
still needed. The bounds leave open rare bases, other budget scales
and collective frequency arguments.
In particular no small-derivative estimate uniform over all triples,
and no advance toward the $abc$ conjecture, follows from these results.

\paragraph{Computational materials.}
The companion for this manuscript is version 0.2.3. Its canonical
repository address, reserved for public distribution, is
\begin{quote}\small\noindent
\href{https://github.com/aconsciousfractal/FCIG-Exact-Fibre-Geometry-Arithmetic-Derivatives}{%
\texttt{https://github.com/aconsciousfractal/}\\
\texttt{FCIG-Exact-Fibre-Geometry-Arithmetic-Derivatives}}
\end{quote}
Anonymous access at this address must be confirmed when the package is released.
The computational evidence is identified by the SHA-256 digest of
\texttt{EVIDENCE\_SHA256.txt}:
\begin{center}\small\texttt{9dba941d261ab87247323478b41bda18dd422b902f4f9f5ad1179d86ccdfe01c}\end{center}
That manifest covers the certificate data, exact checkers, tests and replay
configuration; it excludes the paper and this paragraph, so the identity
does not depend on the PDF. Run
\texttt{python scripts/check\_manifest.py --evidence} to check its entries
and compare the manifest digest with the value printed here.

The accompanying repository has a ten-minute reading path and a
standard-library replay in \texttt{README\_REVIEWER.md} and
\texttt{REPRODUCE.md}. Its verifier checks the $k=47$ factorization,
lattice reduction, optimizer count and primitive fibre; the two
variable-base minima and two finite spectral certificates; and literal
small-cube controls for the prime-family theorem. Mutated certificates
are rejected, and acceptance tests use explicit exceptions so that
Python's optimized mode does not disable them. Such finite checks
support the indicated computations. The universal and density results
rest on the written proofs and the explicitly attributed analytic input.

\paragraph{AI assistance.}
AI tools were used under the author's direction for mathematical
exploration, source comparison, proof criticism, drafting, exact
computations and repository preparation. The accompanying
\texttt{AI\_USE.md} describes their scope. AI cross-review is not
independent human peer review. The human author is responsible for
the mathematics, attribution and any decision to publish.

\appendix
\section{Proof of the common density domain}\label{app:domain}
We prove Theorem~\ref{thm:domain}. All auxiliary bounds are fixed
before the base tends to infinity.

\subsection{The analytic input and its counting consequence}
We use the weighted estimate in Pascadi \cite[Sections 6.1 and 6.3,
Notations 20 and 23, (6.20) and the concluding strict margin]{Pascadi},
with these locators referring to the cited arXiv version.
For a fixed nonnegative nonzero smooth weight at most one on $[1,2]$,
the full log-prime mass has a main term proportional to $X\log X$,
and its truncation at a constant multiple of $X^{1.3}$ misses a fixed
positive proportion. We import that estimate and its analytic and
numerical ingredients. The following counting consequence is a
deduction from the weighted proof, not from bare infinitude.

Each integer $x\le4X$ contributes at most
$\sum_{p\mid x^2+1}\log p\le\log(x^2+1)\le3\log X$ for large $X$.
For fixed $1<\alpha<1.3$, the prime threshold exceeds $(4X)^\alpha$
eventually. Thus some $\eta_0>0$ satisfies, for every sufficiently
large $X$,
\begin{equation}\label{eq:density-input}
 \#\{x\in[X,4X]:P^+(x^2+1)>x^\alpha\}\ge\eta_0 X.
\end{equation}
No restriction on progressions, on prime bases, or on their modular
phase distribution is imported.

\subsection{Small exceptional sets for cosocles and anchors}
For an odd prime $p$, the condition that $p^2$ divide one of
$x,x-1,x+1,x^2+1$ occupies at most five classes modulo $p^2$.
The quadratic has at most two roots, each with a unique lift.
For $x\le Z$, such a prime is at most $Z+1$; therefore the number
of bases with a repeated prime $p>T$ in any of the four values is at most
\begin{equation}\label{eq:square-tail}
 5Z\sum_{m>T}m^{-2}+5\pi(Z+1)\le5Z/T+5\pi(Z+1).
\end{equation}
Here $\pi(Z)=o(Z)$ suffices. One elementary proof fixes a product $Q$
of small primes: apart from finitely many exceptions primes are coprime
to $Q$, so the upper density is at most $\varphi(Q)/Q$.
This tends to zero, since the inverse Euler product contains growing
partial harmonic sums. No prime number theorem is needed.

For fixed $T$, exclusion of divisibility of any of the four values by
$p^{E+1}$ for $p\le T$ has upper density at most
$\sum_{p\le T}5/p^{E+1}$. This is also a valid bound at $2$ when
$E\ge1$, since the quadratic has no root modulo four.
Given $\varepsilon>0$, choose $T$ and then $E$ to make the sum
of these two upper-density bounds less than $\varepsilon$.
Outside the exclusions each cosocle is bounded by
$B_0=\prod_{p\le T}p^{E-1}$, and so $H_0\le B_0^3$.
The degree-two upper bound on $p$ in \eqref{eq:square-tail} is
essential; this argument is not applied to the growing-degree $b$.

For fixed $Y$, CRT gives exact natural densities for the absence
of small simple primes in $x$, and in $x^2+1$ with $p\equiv1\pmod4$:
\begin{equation}\label{eq:anchor-density}
 a(Y)=\prod_{p\le Y}\left(1-\frac{p-1}{p^2}\right),\qquad
 b(Y)=\prod_{\substack{p\le Y\\p\equiv1\ (4)}}
                 \left(1-\frac{2(p-1)}{p^2}\right).
\end{equation}
For the second formula, above each of the two roots modulo $p$
there are $p-1$ lifts of valuation exactly one. Both products tend
to zero. For clarity the needed divergence in the class $1\pmod4$
can be proved by the Euler products, for $s>1$:
\[
 \log\zeta(s)+\log L(s,\chi_4)
       =2\sum_{p\equiv1\ (4)}p^{-s}+O(1),\qquad s\downarrow1.
\]
The higher prime powers contribute a bounded term. The alternating
series $L(s,\chi_4)=1-3^{-s}+5^{-s}-\cdots$ converges uniformly for
$s\ge1$ and lies between $2/3$ and $1$, so its logarithm is bounded.
The divergence of $\zeta(s)$ proves $\sum_{p\equiv1\ (4)}1/p=\infty$.
Use $\log(1-t)\le-t$ in \eqref{eq:anchor-density}.
Independence of the two anchor events is not needed.

\subsection{All large primes paid by reciprocal mass}
We claim, for integer $X\ge2$,
\begin{equation}\label{eq:reciprocal-mean}
 \sum_{x=2}^X\sigma_b<20X,\qquad
 \sum_{x=2}^X(\sigma_x+\sigma_b)<21X.
\end{equation}
For an odd prime $p\le X$, put
\[
 R_p=\rad_{\rm odd}(p-1),\quad
 c_p=\begin{cases}2,&p\equiv3\ (4),\\4,&p\equiv1\ (4),\end{cases}
 \quad m_p=c_pR_p,\quad
 \delta_p=\frac{c_p}{p}\prod_{\ell\mid R_p}(2-1/\ell).
\]
The condition $p\mid b$ says that $d=\ord_p(x)$ divides $N(x)$.
It is equivalent to $d=2^j r\mid m_p$ with $r$ odd squarefree and
$r\mid x+1$. For each such $d$, the $\varphi(d)$ order-$d$ classes
modulo $p$ combine with $x=-1\pmod r$ into as many progressions
modulo $pr$. Summing gives $\sum_{d\mid m_p}\varphi(d)=m_p$
progressions and total density
\[
 \frac1p\sum_j\varphi(2^j)\sum_{r\mid R_p}\frac{\varphi(r)}r
   =\delta_p.
\]
Thus $\#\{2\le x\le X:p\mid b\}\le X\delta_p+m_p$.
Since $m_p\mid p-1$ and
$\prod_{\ell\mid R_p}(2-1/\ell)\le2^{\omega(p-1)}
\le2\sqrt{p-1}$, one has $\delta_p/p<8p^{-3/2}$.
The sum of the weighted counts for odd $p\le X$ is less than
$8X\sum_{n\ge3}n^{-3/2}+\pi(X)<13X$; the prime $2$ adds at most
$X/2$. For primes $p>X$, on each individual base
\[
 \sum_{\substack{p\mid b\\p>X}}\frac1p
 \le\frac{\log b}{X\log X}
 <\frac{4(x+1)\log x}{X\log X}\le6.
\]
This proves the first bound. Finally
$\sum_{x=2}^X\sigma_x=\sum_{p\le X}\lfloor X/p\rfloor/p<X$,
giving the second. The estimate includes every prime divisor of $b$;
it asserts no tightness for its cosocle.

\subsection{Simultaneous selection}
Take $\varepsilon=\eta_0/64$. Choose the cosocle bound $K$ so that
its exceptional upper density is less than $\varepsilon$, and then
fix $Y$ with $a(Y)+b(Y)<\varepsilon$.
For all sufficiently large $X$, each of these two exceptional counts
up to $4X$ is less than $8\varepsilon X=\eta_0X/8$.
Choose $L>336/\eta_0$. Markov's inequality and
\eqref{eq:reciprocal-mean} bound its exception up to $4X$ by
$84X/L<\eta_0X/4$. Subtracting their union from
\eqref{eq:density-input} leaves at least $\eta_0X/2$ bases in
$[X,4X]$. Taking $X$ near $Z/4$ proves positive lower natural
density of the globally defined selected set.

Its simple quadratic anchor $p_b\equiv1\pmod4$ has order four,
does not divide $N$, and stays simple in $b$ by the valuation formula.
The same holds for $P>x^\alpha>x$. Eventually $P>Y$ so it is
different from the small anchor. The source has its required simple
anchor by construction, proving every assertion of
Theorem~\ref{thm:domain}.

\section{A full-domain bound in the fixed-power family}\label{app:fixed}
The following necessary condition helps separate the fixed-power open
problem from the prime-family theorem.

\begin{theorem}\label{thm:fixed}
For $k\ge1$, let $P_1$ be the largest prime of exponent exactly one
in $c=3^k+2$, with $P_1=0$ if there is none. Then
\begin{equation}\label{eq:fixed-bound}
 1\le\frac{\mpr(2,3^k)}{m(2,3^k)}<32\max(5,P_1)^3.
\end{equation}
Unbounded ratios in this family would consequently imply that at
least one of the fixed bases $2,3$ has infinitely many non-Wieferich
primes. This is a necessary consequence, not a converse.
\end{theorem}
\begin{proof}
Take $k\ge2$ and write
\[
 \begin{gathered}
 T=3^{k-1},\quad e_p=v_p(c),\quad A_p=ce_p/p,\quad h_c=c/\rad(c),\\
 g=\gcd_p A_p,\quad\ell=\lcm_p e_p,\quad h=\gcd(kc,3g),\quad t=g/h_c,
 \end{gathered}
\]
\[
 \lambda=\sum_{p\mid c}e_p/p,\quad a_0=1+2\lambda,\quad
 \nu=\frac{hT}{ca_0}.
\]
Eliminating $z_2$ from $F$ gives
$W/T=kc z_3-3\sum A_pz_p$, so $\Dgen=Th$.
Also $h_c\mid h$ and $t\mid\ell$: at a prime dividing $\rad(c)$
inspect its own coefficient $e_p\rad(c)/p$, and at any other prime
inspect the minimum valuation among the $e_p$.
The identity $W=2\sum A_pz_p-cz_2$ on $F=0$ gives $m\ge\nu$.
A real primitive lift is
\[
 z_2^*=-\nu,\quad z_p^*=\nu\ (p\mid c),\quad
 y^*=z_3^*=\frac{(c\lambda+1)\nu}{kT}.
\]
It has norm $\nu$: all primes of $c$ are at least five,
$c<5^k$ gives $\lambda<k/5$, and
$k(2c-10)\ge24>15$ for $c\ge11$ proves $c\lambda+1\le kT$.

Choose $y$ in the solution class $kcy=h\pmod{3g}$ nearest $y^*$.
Its period is $3g/h$, hence $|y-y^*|\le3g/(2h)$.
Then $z_2=(2ky-h)/3$ is integral and
$U=(kcy-h)/3\in g\Z$.
The real $c$-row with all coordinates
$v=\nu+k(y-y^*)/(3\lambda)$ has target $U$.
Put $P_p=p\ell/e_p$ and $Q=c\ell$, so $A_pP_p=Q$.
Start from any integer lift of $U$, keep its residue classes modulo
$P_p$, and round $v$ down in each class.
The deficit divided by $Q$ is an integer in
$\{0,\ldots,\omega(c)-1\}$, since it is a sum of fractional parts.
There are more positive fractional parts than the integer deficit unless it is zero.
Rounding that many coordinates with positive fractional part up once gives an exact lift with
$|z_p-v|<P_p$ in every coordinate. Thus, with $P_{\max}=\max P_p$,
\[
 \mpr\le\nu+\frac{kg}{h}\left(1+\frac1{2\lambda}\right)+P_{\max}.
\]
Dividing by $m\ge\nu$, using $c/T\le5$ and $h\ge h_c$, yields
\begin{equation}\label{eq:fixed-error}
 \frac{\mpr}{m}\le1+5a_0k\frac{t}{h_c}
              \left(1+\frac1{2\lambda}\right)
              +5a_0\frac{P_{\max}}{h_c}.
\end{equation}

Put $M=\max(5,P_1)$, let $B_s$ be the product of the simple primes,
and set $v_0=\log h_c$. Then
\[
 c\le B_sh_c^2,\quad k<3M+2v_0,\quad
 \lambda_s\le M/5,\quad\lambda\le M/5+v_0.
\]
Here $\log B_s\le\sum_{p\le M}\log p<3M$ follows from the dyadic
binomial bound: each interval $(2^{j-1},2^j]$ contributes at most
$\log\binom{2^j}{2^{j-1}}<2^j\log2$, whose sum up to the next
power of two is less than $4M\log2<3M$.
For repeated primes $p\ge5,e\ge2$,
\[
 e\le p^{(e-1)/2},\qquad e/p^{e-1}\le2/5,\qquad e^2/p^e\le4/25.
\]
Each follows at $e=2$ and then by successive ratios. In particular
$t\le\ell\le\sqrt{h_c}$. The elementary maxima
$v_0e^{-v_0/2}<1$ and $v_0^2e^{-v_0/2}<3$ give
\begin{equation}\label{eq:prod-one}
 a_0k/\sqrt{h_c}<\tfrac65M^2+\tfrac{49}5M+14\le4M^2.
\end{equation}
We also have
\begin{equation}\label{eq:prod-two}
 a_0P_{\max}/h_c<M^2.
\end{equation}
To see this, write $f_{p,e}=e/p^{e-1}$.
For a selected simple $p$, $P_p/h_c\le M\prod_{q\text{ repeated}}f_q$
and, if there are $n$ repeated primes,
$\lambda_{\rm rep}P_p/h_c\le(4M/25)n(2/5)^{n-1}\le4M/25$
(the sum is zero for $n=0$). Thus
$a_0P_p/h_c\le M(1+2M/5+8/25)<M^2$.
For a repeated selected $p$, $P_p/h_c\le1$; its own contribution
to $\lambda_{\rm rep}P_p/h_c$ is at most $2/5$, and all other
ones together at most $4/25$. Hence
$a_0P_p/h_c\le1+2M/5+28/25<M^2$, proving \eqref{eq:prod-two}.

If $B_s>1$, $\lambda\ge1/M$, and \eqref{eq:prod-one} bounds the
first error in \eqref{eq:fixed-error} by
$20M^2(1+M/2)\le14M^3$. Together with \eqref{eq:prod-two} this
gives $\mpr/m<16M^3$.
If $B_s=1$, $c=2\pmod3$ is not a square, so some repeated
exponent $e_p$ is odd and at least three. Since $\lambda\ge e_p/p$,
\[
 \frac{\ell}{h_c\lambda}
 \le p^{2-e_p}\prod_{q\ne p}\frac{e_q}{q^{e_q-1}}
 \le h_c^{-1/2}.
\]
Now $k<2v_0$, $a_0\le1+2v_0$, and the first error in
\eqref{eq:fixed-error} is less than
$(15/2)(2v_0+4v_0^2)e^{-v_0/2}<105$.
With $M=5$, \eqref{eq:prod-two} gives $\mpr/m<231<32M^3$.
For $k=1$, direct substitution gives $m=1,\mpr=2$.

Finally, for $p\ge5$ dividing $3^k+2$, put
$f_p(a)=(a^{p-1}-1)/p\pmod p$.
Raising $3^k=-2+p((3^k+2)/p)$ to $p-1$ modulo $p^2$ yields
\[
 (3^k+2)/p\equiv2(k f_p(3)-f_p(2))\pmod p.
\]
If $p$ is simple, the left side is nonzero, so $p$ is non-Wieferich
to at least one of $2,3$. An unbounded ratio forces unbounded $P_1$
by \eqref{eq:fixed-bound}; a union of two finite sets is finite.
This proves the fixed-base alternative.
\end{proof}

\section{Weighted covering of the type-\texorpdfstring{\(A\)}{A} lattice}\label{app:radius}
Let \(r\ge2\), \(q_1,\ldots,q_r>0\), and
\[
 H=\{y\in\R^r:\textstyle\sum_i y_i=0\},\quad
 A_{r-1}=H\cap\Z^r,\quad N_q(y)=\max_i q_i|y_i|.
\]
We retain the following as a useful geometric proposition, with its metric
explicit. It is not required by the minimum certificate in Section~\ref{sec:k47}.

\begin{proposition}\label{prop:radius}
The covering radius is
\begin{equation}\label{eq:radius}
 \mu(A_{r-1},N_q)=
 \max\left\{\frac12\max_iq_i,\,
                 \frac{r-1}{\sum_i1/q_i}\right\}.
\end{equation}
For \(r=1\) the radius is zero instead.
\end{proposition}
\begin{proof}
Let \(\mu_*\) be the right side. For \(y\in H\), set
\[
 l_i=\lceil y_i-\mu_*/q_i\rceil,\qquad
 u_i=\lfloor y_i+\mu_*/q_i\rfloor.
\]
Each interval contains an integer, since \(2\mu_*/q_i\ge1\).
The floor inequality gives
\[
 \sum_i u_i\ge
 \left\lfloor\mu_*\sum_i1/q_i\right\rfloor-(r-1)\ge0.
\]
Applied to \(-y\), it gives \(\sum_i l_i\le0\).
Sums of integer intervals have no holes, so integers
\(z_i\in[l_i,u_i]\) can be chosen with sum zero. This proves the upper bound.

For a maximal-weight index \(i\) and \(j\ne i\), take
\(y_i=1/2,y_j=-1/2\) and all other coordinates zero.
Its distance to every integer vector is at least \(q_i/2\).
For the other bound put \(R_0=(r-1)/\sum_i1/q_i\), and choose
\(y_1=-R_0/q_1\), \(y_i=1-R_0/q_i\) for \(i>1\).
This point lies in \(H\). If \(N_q(y-z)<R_0\), then
\(z_1<0\) and \(z_i<1\) for \(i>1\).
For integer \(z\) this implies \(\sum_i z_i\le-1\), impossible.
Both lower bounds follow.
\end{proof}

For pairwise coprime integers \(q_i\ge2\), put \(Q=\prod_iq_i\) and
\(w_i=Q/q_i\). Modular reduction gives
\[
 \ker_{\Z}w=\operatorname{diag}(q_i)A_{r-1}.
\]
The diagonal map transports the original maximum norm to \(N_q\).
Since \(\gcd_i w_i=1\), every integer target \(t\) has a lift.
Its real minimum is \(\rho(t)=|t|/\sum_i w_i\).
If \(m_w(t)=\min_{w x=t,\ x\in\Z^r}\norm{x}\), then
\begin{equation}\label{eq:lift-radius}
 \lceil\rho(t)\rceil\le m_w(t)\le
 \lfloor\rho(t)+\mu(A_{r-1},N_q)\rfloor.
\end{equation}
For the upper bound, subtract any integer lift from a real minimizer,
approximate the resulting point of the kernel by a kernel lattice point,
and translate back. The norm of an integer vector is integral.

Prescribed-sum rounding is studied by Cont and Heidari
\cite[\S\S3.1--3.5]{ContHeidari}, including common or symmetric losses.
Unequal coordinate weights in a maximum norm do not give a symmetric
loss under coordinate permutations, so that comparison is not an
automatic import of \eqref{eq:radius}.
No priority claim for the formula is made here.
Furthermore, a covering radius controls distance from arbitrary points
of the real kernel. It cannot be substituted for a radial lifting
overhead, or assigned to a projected fibre, without identifying both
the actual difference lattice and the transported norm.

\begin{thebibliography}{10}\raggedright
\bibitem{Pasten}
H. Pasten, \emph{Arithmetic derivatives through geometry of numbers},
Canadian Mathematical Bulletin \textbf{65} (2022), no.~4, 906--923.
\href{https://doi.org/10.4153/S0008439521000990}{doi:10.4153/S0008439521000990}.
Consulted \href{https://arxiv.org/abs/2106.16165v3}{arXiv:2106.16165v3}
(11 December 2021).

\bibitem{Stanley}
R. P. Stanley, \emph{Smith normal form in combinatorics},
Journal of Combinatorial Theory, Series A \textbf{144} (2016), 476--495.
\href{https://doi.org/10.1016/j.jcta.2016.06.013}{doi:10.1016/j.jcta.2016.06.013}.
Consulted \href{https://arxiv.org/abs/1602.00166v2}{arXiv:1602.00166v2}.

\bibitem{UniqueBetti}
P. A. Garc\'ia-S\'anchez, I. Ojeda and J. C. Rosales,
\emph{Affine semigroups having a unique Betti element},
Journal of Algebra and Its Applications \textbf{12} (2013), no.~3,
1250177, 11 pp.
\href{https://doi.org/10.1142/S0219498812501770}{doi:10.1142/S0219498812501770}.
Consulted \href{https://arxiv.org/abs/1203.4138v2}{arXiv:1203.4138v2}.

\bibitem{Autry}
J. Autry, P. Graves, J. Loucks, C. O'Neill, V. Ponomarenko and S. Yih,
\emph{Squarefree divisor complexes of certain numerical semigroup elements},
Involve \textbf{14} (2021), no.~1, 1--9.
\href{https://doi.org/10.2140/involve.2021.14.1}{doi:10.2140/involve.2021.14.1}.
Consulted author text: \url{https://vadim.sdsu.edu/aglopy.pdf}.

\bibitem{CyrusianKaplan}
S. Cyrusian and N. Kaplan, \emph{Ordinarization numbers of numerical semigroups},
Communications in Algebra (2026), 1--24.
\href{https://doi.org/10.1080/00927872.2026.2615092}{doi:10.1080/00927872.2026.2615092}.
Consulted \href{https://arxiv.org/abs/2506.10222v2}{arXiv:2506.10222v2}.

\bibitem{ContHeidari}
R. Cont and M. Heidari, \emph{Optimal rounding under integer constraints},
\href{https://arxiv.org/abs/1501.00014v3}{arXiv:1501.00014v3}
(5 August 2026; first version 2014).

\bibitem{Pascadi}
A. Pascadi, \emph{Large sieve inequalities for exceptional Maass forms
and the greatest prime factor of $n^2+1$},
Forum of Mathematics, Pi \textbf{14} (2026), e8.
\href{https://doi.org/10.1017/fmp.2026.10025}{doi:10.1017/fmp.2026.10025}.
All numbered locators here refer to
\href{https://arxiv.org/abs/2404.04239v3}{arXiv:2404.04239v3}.

\bibitem{FelliniMurty}
N. Fellini and M. R. Murty,
\emph{Wieferich primes in number fields and the conjectures of
Ankeny--Artin--Chowla and Mordell},
Journal of Number Theory \textbf{285} (2026), 209--229.
\href{https://doi.org/10.1016/j.jnt.2026.01.002}{doi:10.1016/j.jnt.2026.01.002}.
Theorem locators refer to
\href{https://arxiv.org/abs/2508.08472v2}{arXiv:2508.08472v2}.
\end{thebibliography}
\end{document}
